%% start of file `template.tex'.
%% Copyright 2006-2013 Xavier Danaux (xdanaux@gmail.com).
%
% This work may be distributed and/or modified under the
% conditions of the LaTeX Project Public License version 1.3c,
% available at http://www.latex-project.org/lppl/.


\documentclass[11pt,a4paper,roman]{moderncv}        % possible options include font size ('10pt', '11pt' and '12pt'), paper size ('a4paper', 'letterpaper', 'a5paper', 'legalpaper', 'executivepaper' and 'landscape') and font family ('sans' and 'roman')
%\usepackage{fontawesome} % Pacote para os ícones
% modern themes
\moderncvstyle{banking}   
% style options are 'casual' (default), 'classic', 'oldstyle' and 'banking'
\moderncvcolor{blue}                                % color options 'blue' (default), 'orange', 'green', 'red', 'purple', 'grey' and 'black'
%\renewcommand{\familydefault}{\sfdefault}         % to set the default font; use '\sfdefault' for the default sans serif font, '\rmdefault' for the default roman one, or any tex font name
\nopagenumbers{}                                  % uncomment to suppress automatic page numbering for CVs longer than one page

% character encoding
\usepackage[utf8]{inputenc}
\usepackage{fontawesome5}
\usepackage{fontspec}
\usepackage{tabularx}
\usepackage{ragged2e}
% if you are not using xelatex ou lualatex, replace by the encoding you are using
%\usepackage{CJKutf8}                              % if you need to use CJK to typeset your resume in Chinese, Japanese or Korean

% adjust the page margins
\usepackage[scale=0.8]{geometry}
\usepackage{multicol}
%\setlength{\hintscolumnwidth}{3cm}                % if you want to change the width of the column with the dates
%\setlength{\makecvtitlenamewidth}{10cm}           % for the 'classic' style, if you want to force the width allocated to your name and avoid line breaks. be careful though, the length is normally calculated to avoid any overlap with your personal info; use this at your own typographical risks...

\usepackage{import}

% personal data
\name{Leonardo}{Santos}
% \title{Curriculum Vitae}                               % optional, remove / comment the line if not wanted
%\address{Pinhais, Paraná,}{}{}% optional, remove / comment the line if not wanted; the "postcode city" and and "country" arguments can be omitted or provided empty
% \phone[mobile]{909-839-3097}                   % optional, remove / comment the line if not wanted
% \phone[fixed]{01234 123456}                    % optional, remove / comment the line if not wanted
%\phone[fax]{+3~(456)~789~012}                      % optional, remove / comment the line if not wanted
% \email{xpan1@swarthmore.edu}                               % optional, remove / comment the line if not wanted
% \homepage{shawnpan.me}                         % optional, remove / comment the line if not wanted
% \extrainfo{}                 % optional, remove / comment the line if not wanted
%\photo[64pt][0.4pt]{picture}                       % optional, remove / comment the line if not wanted; '64pt' is the height the picture must be resized to, 0.4pt is the thickness of the frame around it (put it to 0pt for no frame) and 'picture' is the name of the picture file
%\quote{Some quote}                                 % optional, remove / comment the line if not wanted

% to show numerical labels in the bibliography (default is to show no labels); only useful if you make citations in your resume
%\makeatletter
%\renewcommand*{\bibliographyitemlabel}{\@biblabel{\arabic{enumiv}}}
%\makeatother
%\renewcommand*{\bibliographyitemlabel}{[\arabic{enumiv}]}% CONSIDER REPLACING THE ABOVE BY THIS

% bibliography with mutiple entries
%\usepackage{multibib}
%\newcites{book,misc}{{Books},{Others}}
  
\newcommand*{\customcventry}[7][.25em]{
  \begin{tabular}{@{}l} 
    {\bfseries #4}
  \end{tabular}
  \hfill% move it to the right
  \begin{tabular}{l@{}}
     {\bfseries #5}
  \end{tabular} \\
  \begin{tabular}{@{}l} 
    {\itshape #3}
  \end{tabular}
  \hfill% move it to the right
  \begin{tabular}{l@{}}
     {\itshape #2}
  \end{tabular}
  \ifx&#7&%
  \else{\\%
    \begin{minipage}{\maincolumnwidth}%
      \small#7%
    \end{minipage}}\fi%
  \par\addvspace{#1}}

\newcommand*{\customcvproject}[4][.25em]{
%   \vfill\noindent
  \begin{tabular}{@{}l} 
    {\bfseries #2}
  \end{tabular}
  \hfill% move it to the right
  \begin{tabular}{l@{}}
     {\itshape #3}
  \end{tabular}
  \ifx&#4&%
  \else{\\%
    \begin{minipage}{\maincolumnwidth}%
      \small#4%
    \end{minipage}}\fi%
  \par\addvspace{#1}}

\setlength{\tabcolsep}{12pt}

%----------------------------------------------------------------------------------
%            content
%----------------------------------------------------------------------------------
\begin{document}
%\begin{CJK*}{UTF8}{gbsn}                          % to typeset your resume in Chinese using CJK
%-----       resume       ---------------------------------------------------------
\makecvtitle
\vspace*{-10mm}

\begin{center}
	\makebox[\textwidth][c]{ % Faz a tabela ocupar a largura total da página e centraliza
		\begin{tabular}{c c c c}
			\faLinkedin\enspace linkedin.com/in/leonardo-20-as & 
			\faEnvelope\enspace leonardo\_as20@hotmail.com & 
			\faGithub\enspace Je-Leo-AS & 
			\faMobile\enspace (41) 988267727 \\  
		\end{tabular}
	}
\end{center}


\section{Resumo}
{\customcventry{}{}{Engenheiro eletricista | Cientista de Dados | Desenvolvedor de Software | Desenvolvedor de Firmware}{}{}{\parbox{\textwidth}{Engenheiro eletricista com experiência em projetos elétricos, design de circuitos integrados, firmware e software. Atuei na pesquisa e desenvolvimento de sistemas digitais para processamento de sinais, incluindo a modelagem e implementação de um Pré-Distorcedor Digital (DPD) para comunicação sem fio e geradores de números aleatórios em FPGA. Durante esse período, desenvolvi modelos matemáticos para correção de não linearidades em amplificadores de potência, o que fortaleceu minha atuação como cientista de dados ao aplicar técnicas de modelagem e otimização para resolver problemas complexos. Também fui monitor de Eletrônica Digital e Microeletrônica e participei de competições e projetos de extensão. Atualmente, trabalho com desenvolvimento de software para agricultura de precisão, focado no armazenamento e processamento de grandes bases de dados de equipamentos agrícolas. Tenho experiência com Python, VHDL, FPGA, C++, Cadence, Xilinx Vivado e machine learning, buscando sempre inovação e eficiência nos projetos em que atuo.}}}


\section{Formação acadêmica}
{\customcventry{Março 2019 - Dezembro 2024}{Discente de Engenharia elétrica}{Universidade federal do Paraná}{Curitiba, PR}{}{}{\parbox{\textwidth}{Durante minha graduação, atuei como monitor nas disciplinas de Eletrônica Digital e Microeletrônica, auxiliando alunos na resolução de exercícios, esclarecimento de dúvidas e desenvolvimento de práticas de laboratório. Também ministrei um curso de Python para estudantes de Engenharia Química e fui gerente administrativo da Semana de Atualização Acadêmica de Engenharia Elétrica (Seatel), onde organizei cursos, palestras e visitas técnicas.	
No Programa de Educação Tutorial (PET), participei de projetos como IoPET (focado em IoT e firmware), BigPET (machine learning, onde fui líder) e PET3D (design e impressão 3D, também como líder). 			
Integrei um projeto de extensão para o desenvolvimento de um veículo off-road, atuando no subsistema elétrico, onde desenvolvi e validei o chicote elétrico do veículo, além de trabalhar com design de PCBs, programação de firmware e modelagem 3D no software Catia. Com a equipe, conquistamos: 
\begin{itemize}
	\item 1º lugar no projeto elétrico no Regional Sul 2019;
	\item 11º lugar nacional no Baja SAE 2020;
	\item 2º lugar no projeto elétrico no Regional Sul 2020.
\end{itemize}}}}

\subsection{}
{\customcventry{Março 2025 - Atualmente}{Mestrando em Engenharia Elétrica}{Universidade Federal do Paraná}{Curitiba, PR}{}{}{\parbox{\textwidth}{Mestrando na área de desenvolvimento de circuitos digitais para sistemas de telecomunicações, com foco em técnicas avançadas de processamento de sinais e otimização de hardware. Pesquiso arquiteturas eficientes para FPGA e ASIC, visando melhorar o desempenho e a eficiência energética de sistemas embarcados em aplicações de comunicação sem fio.}}}


\section{Atividades complementares}
{\customcventry{Janeiro 2019 - 2021}{Membro do subsistema da elétrica e gestão}{UFPR BAJA SAE}{Curitiba, PR}{}{}{\parbox{\textwidth}{Participei de um projeto de extensão focado no desenvolvimento de um veículo off-road para aplicar, na prática, o conhecimento adquirido em sala de aula. O projeto era dividido em subsistemas, cada um responsável por uma etapa do desenvolvimento do carro. Atuei no subsistema elétrico, responsável pelos testes e validação dos demais subsistemas, além do projeto do chicote elétrico, representando o sistema elétrico do veículo. Nesse projeto, tive meu primeiro contato com eletrônica prática, envolvendo design de PCBs e programação de firmware. Também desenvolvi habilidades em modelagem 3D usando o software Catia para o projeto do chicote elétrico. Durante o tempo no projeto, alcançamos as seguintes colocações nas competições Baja SAE: Regional Sul 2019: 1º lugar de projeto elétrico; Nacional 2020: 11º lugar de projeto elétrico; Regional Sul 2020: 2º lugar de projeto elétrico Na competição Baja SAE BRASIL, os alunos formam equipes para representar sua Instituição de Ensino Superior, participando de avaliações anuais comparativas de seus projetos. No Brasil, a etapa nacional é chamada Competição Baja SAE BRASIL, e as regionais incluem as etapas Sul, Sudeste e Nordeste.}}}

\subsection{}
{\customcventry{Agosto 2021 - Dezembro 2024}{Pesquisador científico}{GICS}{Curitiba, PR}{}{}{\parbox{\textwidth}{Entre agosto de 2021 e agosto de 2022, tive meu primeiro contato com VHDL, trabalhando no desenvolvimento de um circuito integrado baseado em um gerador de números verdadeiramente aleatórios (TRNG), partindo de um projeto previamente elaborado. O cronograma incluiu o estudo de VHDL e dispositivos lógicos programáveis (FPGAs), seguido da implementação e validação da arquitetura do TRNG em FPGA. Após sua validação, o design foi implementado como circuito integrado utilizando as ferramentas da Cadence. De 2022 a 2023, atuei na modelagem de um pré-distorcedor digital baseado no Polinômio de Memória, empregando Python para otimizar sistemas de processamento digital de sinais, trabalho que foi posteriormente desenvolvido em VHDL. Entre 2023 e 2024, realizei a implementação em hardware desse sistema, inicialmente utilizando FPGAs com as ferramentas da Xilinx. Após a validação, o design do circuito integrado foi desenvolvido com as ferramentas da Cadence.}}}

\subsection{}
{\customcventry{Agosto 2021 - Dezembro 2024}{Gerente de projeto | Integrante}{Projeto Extensão Tutorial}{Curitiba, PR}{}{}
{\parbox{\textwidth}{O Programa de Educação Tutorial (PET) tem como compromisso fundamental aprimorar os cursos de graduação. Durante minha participação no programa, estive envolvido ativamente nos seguintes projetos:
\begin{itemize}
	\item IoPET: Projeto focado no desenvolvimento de atividades relacionadas à Internet das Coisas (IoT), programação de software e firmware. Contribuí com o desenvolvimento de uma fechadura eletrônica utilizando o módulo ESP01, que se comunicava por meio do protocolo MQTT.
	\item BigPET: Projeto voltado ao estudo de machine learning e ao desenvolvimento de propostas e desafios na área. Atuei como líder do projeto, ministrei um curso de introdução à machine learning para a turma de Engenharia Química e participei do desenvolvimento de um detector de fake news utilizando uma biblioteca de vetorização de strings.
	\item PET3D: Projeto dedicado à familiarização com ferramentas de design, modelagem e impressão 3D. Também como líder do projeto, desenvolvemos suportes para máscaras face shield durante a pandemia, em parceria com hospitais, como forma de contribuir no combate à COVID-19 e oferecer suporte às equipes de saúde.
\end{itemize}
Além dessas atividades, organizei cursos durante a semana acadêmica e atuei como monitor nas disciplinas de Eletrônica Digital e Microeletrônica, promovendo a integração entre teoria e prática para os estudantes.}}}

\section{Experiência}

{\customcventry{Outubro de 2022 - Atualmente}{Desenvolvedor de software / Cientista de Dados}{IoTag Tecnologia}{Curitiba, PR}{}
{Na empresa, nosso objetivo é desenvolver soluções inovadoras em agricultura de precisão, com foco em firmware, algoritmos de machine learning e sistemas embarcados, visando aumentar a eficiência e a automação no setor agrícola. Participei ativamente da implementação de firmware em C e C++ para a criação da interface do terminal virtual de maquinário agrícola, além de desenvolver algoritmos de ciência de dados aplicados a dados espaciais e à coleta de dados dos equipamentos. Também criei serviços que integram bancos de dados e plataformas de cloud computing, como AWS, garantindo a escalabilidade e a eficiência das soluções desenvolvidas.}

\subsection{}

{\customcventry{Junho 2022 – Outubro 2022}{Estagiário de engenharia}{Lactec}{Curitiba, PR}{}
{Durante o estágio, auxiliei na instrumentação do projeto de desenvolvimento de um gerador de energia ondomotriz. Colaborei no desenvolvimento de sensores, como o sensor de nível de água, utilizando tecnologia capacitiva para medições precisas, e o sensor de fluxo, implementado com um motor elétrico, para monitorar e otimizar o desempenho do sistema. Essa experiência me permitiu aprimorar habilidades em eletrônica, programação de microcontroladores e integração de sistemas, contribuindo significativamente para o avanço do projeto.
}

\section{Projetos}

{\customcvproject{Projeto de assistente virtual}{}
  {Este projeto é um assistente virtual multimodal, capaz de interpretar comandos de voz e texto e responder de maneira inteligente e contextualizada. Ele utiliza o LangChain para processamento de linguagem natural e o FastAPI para gerenciar o backend, garantindo respostas rápidas e eficientes. Além disso, conta com um sistema de reconhecimento de fala que permite interações naturais, possibilitando que o usuário se comunique por voz e receba respostas tanto em áudio quanto em texto. O frontend da aplicação foi desenvolvido em JavaScript, proporcionando uma interface simples e intuitiva para facilitar a interação com o assistente. Esse chatbot pode ser utilizado para diversas finalidades, como atendimento ao cliente, assistentes pessoais e integração com sistemas inteligentes, tornando a experiência do usuário mais fluida e dinâmica.
  }
}
\subsection{}
{\customcvproject{Fechadura eletrônica}{}{
Este projeto teve como objetivo desenvolver uma fechadura eletrônica inteligente utilizando um ESP-01 (módulo ESP8266) para receber comandos via protocolo MQTT e acionar uma fechadura eletromagnética do tipo solenóide. A comunicação MQTT permite que a fechadura seja controlada remotamente por meio de um aplicativo, sistema web ou até mesmo integração com assistentes de voz e automação residencial.
}
\subsection{}
{\customcvproject{Projeto de um circuito integrado de um TRNG}{}
{A evolução da tecnologia aumentou a demanda por sistemas rápidos e seguros, impulsionando a pesquisa em criptografia. Números aleatórios são essenciais para algoritmos criptográficos, sendo gerados por PRNGs ou TRNGs. Este projeto desenvolveu um TRNG em tecnologia de 130 nm, baseado em um projeto anterior em VHDL. O circuito foi otimizado, sintetizado com Genus e validado por simulações. Após posicionamento e roteamento no Innovus, obteve-se um circuito com 43 células lógicas e consumo de 41 µW.
}}
\subsection{}
{\customcvproject{Detector de Fakenews}{}
	{Sistema desenvolvido para identificar a veracidade de notícias utilizando técnicas de processamento de texto e aprendizado de máquina. O projeto envolveu a leitura e limpeza de dados, remoção de stopwords, aplicação de lematização e treinamento de modelos como Regressão Logística, KNN e Random Forest.
	}
}

\section{ADICIONAIS}
\begin{minipage}{\maincolumnwidth}%
	\small{
		\begin{itemize}
			\item Cursos concluídos: Desenvolvimento Android, Machine Learning e Banco de Dados (60 horas cada)
			\item Linguagens de Programação: Python, C, C++, Java, JavaScript, PHP, HTML/CSS
			\item Fluente em Português, proficiência avançada em Inglês
			\item Habilidades Técnicas:
			\begin{itemize}
				\item Experiência em modelagem 3D com SolidWorks, Fusion 360 e Catia.
				\item Projeto e design de circuitos integrados utilizando as ferramentas Cadence e ADS.
				\item Simulação de circuitos eletrônicos com Proteus.
				\item Projeto de PCB com Altium Designer.
				\item Desenvolvimento em Python com frameworks como OpenCV, FastAPI e LangChain.
				\item Experiência com bancos de dados e SQL.
				\item Desenvolvimento web.
				\item Programação em Java.
				\item Boa comunicação interpessoal.
				\item Capacidade de planejamento eficiente.
				\item Gestão eficaz do tempo.
			\end{itemize}
	\end{itemize}}%
\end{minipage}%

}
% Publications from a BibTeX file without multibib
%  for numerical labels: \renewcommand{\bibliographyitemlabel}{\@biblabel{\arabic{enumiv}}}% CONSIDER MERGING WITH PREAMBLE PART
%  to redefine the heading string ("Publications"): \renewcommand{\refname}{Articles}
%\nocite{*}
%\bibliographystyle{plain}
%\bibliography{publications}                        % 'publications' is the name of a BibTeX file

% Publications from a BibTeX file using the multibib package
%\section{Publications}
%\nocitebook{book1,book2}
%\bibliographystylebook{plain}
%\bibliographybook{publications}                   % 'publications' is the name of a BibTeX file
%\nocitemisc{misc1,misc2,misc3}
%\bibliographystylemisc{plain}
%\bibliographymisc{publications}                   % 'publications' is the name of a BibTeX file

%-----       letter       ---------------------------------------------------------

\end{document}


%% end of file `template.tex'.
